\subsection{DOMINANT CHARACTERS}

Denote by \(X_{+}\) the set of elements \(\lambda\) of \(\mathrm{X(T)}\) such that \(\langle\lambda,x\rangle\geq0\) for all \(x\in\Gamma(\mathrm{T})_{++}\) , that is, such that the linear form \(\delta(\lambda):\mathfrak{t}_{\mathbf{C}}\to\mathbf{C}\) maps the chamber C of t to \(iR_{+}\) .

Give X(T) the ordered group structure for which the positive elements are those of \(X_{+}\) ; put \(R_{+} = R(G, T) \cap X_{+}\) and \(R_{-} = -R_{+}\) . The elements of \(R_{+}\) are called positive roots, those of \(R_{-}\) negative roots; every root is either positive or negative (Chap. VI, §1, no. 6, Th. 3). A positive root that is not the sum of two positive roots is said to be simple; every positive root is a sum of simple roots (loc. cit.); the simple roots form a basis of the subgroup of X(T) generated by the roots, a subgroup that can be identified with X(T/C(G)) (§4, no. 4); the reflections with respect to the simple roots generate the Weyl group \(W = W_{G}(T)\) (Chap. VI, §1, no. 5, Th. 2).

Lemma 1. Let \(\lambda\) be an element of \(\mathrm{X(T)}\) . The following conditions are equivalent:

\begin{enumerate}
\def\labelenumi{(\roman{enumi})}
\item
  \(\lambda - w(\lambda) \geq 0\) (resp. \(>0\) ) for all \(w \in \mathrm{W}\) such that \(w \neq 1\) ;
\item
  for all \(w \in \mathrm{W}\) such that \(w \neq 1\) , \(\lambda - w(\lambda)\) is a linear combination with positive coefficients (resp. positive coefficients not all zero) of simple roots;
\item
  \(\langle \lambda, K_{\alpha} \rangle \geq 0\) (resp. \(>0\) ) for every positive root \(\alpha\) ;
\item
  \(\langle\lambda,K_{\alpha}\rangle\geq0\) (resp. >0) for every simple root \(\alpha\) .
\end{enumerate}

The equivalence if (iii) and (iv) is immediate. Since the set of the \(K_{\alpha}\) can be identified with the inverse root system of \(\mathrm{R}(\mathrm{G},\mathrm{T})\) (§4, no. 5), the equivalence of (i) and (iii) follows from Chap. VI, §1, no. 6, Prop. 18 and Cor. The implication (ii) \(\Rightarrow\) (i) is trivial, and the opposite implication follows from loc. cit.

Denote by \(X_{++}\) the set of elements of \(\mathrm{X(T)}\) such that \(\langle\lambda,K_{\alpha}\rangle\geq0\) for every positive root \(\alpha\) . The elements of \(X_{++}\) are said to be dominant. They form a fundamental domain for the operation of W on \(\mathrm{X(T)}\) (Chap. VI, §1, no. 10). We have \(X_{++}\subset X_{+}\) .

If G is simply-connected, for each simple root \(\alpha\) there exists an element \(\varpi_{\alpha}\) of X(T) such that \(\langle\varpi_{\beta}, K_{\alpha}\rangle = \delta_{\alpha\beta}\) for every simple root \(\beta\) , that is, \(s_{\alpha}(\varpi_{\alpha}) = \varpi_{\alpha} - \alpha\) , \(s_{\beta}(\varpi_{\alpha}) = \varpi_{\alpha}\) for every simple root \(\beta \neq \alpha\) ; the \(\varpi_{\alpha}\) are called the fundamental dominant weights; they form a basis of the commutative group \(\mathrm{X(T)}\) and of the commutative monoid \(\mathrm{X}_{++}\) ; more precisely, every element \(\lambda\) of \(\mathrm{X(T)}\) can be written in the form \(\lambda = \sum_{\alpha} \langle \lambda, K_{\alpha} \rangle \varpi_{\alpha}\) .

Denote by \(\rho\) the element of \(\mathrm{X(T)}\otimes \mathbf{Q}\) such that

\[
2 \rho = \sum_ {\alpha \in \mathrm{R} _ {+}} \alpha .
\]

We have \(\langle\rho,K_{\alpha}\rangle=1\) for every simple root \(\alpha\) (Chap. VI, §1, no. 10, Prop. 29). If G is simply-connected, \(\rho\) is the sum of the fundamental dominant weights.

